\chapter{Trust Establishment}

\label{chap:trust-establishment}

\begin{abstract}
In cooperative systems the parties make no assumptions about channel security; this stage is therefore called \emph{trust establishment}.  
Classical \emph{unauthenticated} key-exchange schemes provide confidentiality against passive eavesdroppers, yet they fall to active man-in-the-middle attacks.  
Cooperation cannot rely on unverified assumptions.  
Our sole assumption is access to a local true-random generator; parties remain vulnerable only to signal-loss. 

We propose that each participant first isolates and safeguards an experiment to measure a truly random source, then models the resulting noise according to their individual preferences.  
No harm arises if other parties learn one another’s secret keys, secrets are tactical disadvantages. 
\end{abstract}


\section{Topic}
Let $\Lambda$ model $n$ raw entropy sources $\Pi^{}_{n}$, where $\lambda^{i}$ are the \emph{eigenvalues} of the covariance.  

\begin{protocol}
\caption{Private-key derivation from a random signal.}
\label{prot:eczkmh}
\begin{center}
\[
\mcaln{} \coloneqq (-1,1) \subset \mathbb{R}, 
\qquad 
\mcalnp{} =\mcaln{} \: \bmod{} p,
\]
\[
\wkmodel{},
\]
\[
\mcalk{}.
\]
\[
\mathcal N := (-1,1)\subset\mathbb R,
\qquad
\mathcal N_p := Q_p\!\bigl[\mathcal N\bigr]\subset\mathbb Z_p,
\;\;
Q_p(x)=\bigl\lfloor(x+1)\tfrac{p-1}{2}\bigr\rfloor \pmod{p}.
\]
\end{center}
\end{protocol}

\begin{algorithm}[H]
\caption{Entropy-Maximised Private-Key Derivation (\( \mathsf{Alt}\,1 \))}
\label{alg:emp-max}
\KwIn{Raw entropy batch $\Pi_n$, participant index $i$, target prime $p$}
\KwOut{Private key $k_n^{\,i}$}

\textbf{Isolate} local entropy experiment\tcc*[r]{Physical shielding}
\nl $\pi \gets \textsc{MeasureRandomSource}(\Pi_n)$\;
\nl $\Lambda^{i} \gets \textsc{SolveMaxEntropyModel}\bigl(\pi,\mathrm{constraints}_i\bigr)$\;
\nl $\mathcal{N}_p \gets \textsc{OptimiseQuantiser}\bigl(\Lambda^{i},p\bigr)$\;
\nl $u \gets \Lambda^{i}(\pi)$\tcp*{continuous $(0,1)$ value}
\nl $k_n^{\,i} \gets \bigl\lfloor u\cdot p \bigr\rfloor \bmod p$\;
\nl $k_n^{\,i} \gets \textsc{ExtractUniformBits}\!\bigl(k_n^{\,i}\bigr)$\tcp*{e.g.\ HKDF}
\KwRet{$k_n^{\,i}$}
\end{algorithm}


Step~\ref{step:model} of Algorithm~\ref{alg:pkd}
embeds user preferences into~$\Lambda^{i}$.





This is not finish, algorithm is just a sketch. 



\section{Coordination vs Secrecy}

Independent, perfectly random seeds do not let isolated parties land on the same secret value; information-theoretic results forbid it.


In this sense we can either have, secrecy or coordination. 

| Setting                                                                                         | Mutual information $I(X;Y)$ between Alice’s and Bob’s raw bits | Can a shared secret key be distilled?                                                                                                              |
| ----------------------------------------------------------------------------------------------- | -------------------------------------------------------------- | -------------------------------------------------------------------------------------------------------------------------------------------------- |
| **Truly independent RNGs**                                                                      | $I=0$                                                          | **No.** Shannon (1949) already implies the eavesdropper can simulate either party as well as the other.                                            |
| **Correlated observations of the same phenomenon** (noise on a wire, channel reciprocity, etc.) | $I>0$ but $<H$                                                 | **Yes** – Maurer/Wolf key-agreement theorems (1993-2003) show you can squeeze a key of length $\le I$ by error-correction + privacy-amplification. |
| **Pre-shared seed / entanglement / trusted beacon**                                             | Explicit common randomness                                     | **Trivial** – seed → PRF → key stream.                                                                                                             |



When I(X;Y)=0, the lemma’s bound collapses to “no extractor can give both parties matching uniform bits except with guessing probability.”

Note: Market is a shared channel where I(X;Y)>0 

Unique-optimal-strategy coordination
If a publicly-known optimisation problem has a single global optimum, then any rational agent will converge to the same action deterministically.
Example: both routers run Dijkstra on the same published topology file; they pick the same shortest path and forward packets identically.
Limit: an eavesdropper can run the same algorithm, so secrecy ≈ 0.

%%%% ------------------------------
\bibliographystyle{plain}
\bibliography{references}
